
% =====================================================================
\section{Discussion: the status of the principal statements}
\label{sec:status}
% =====================================================================

A paper that contains theorems, numerical results, hypotheses and assumptions side by side owes the
reader a clear account of which is which. Table~\ref{tab:status} lists the principal statements of
this paper with their nature and the basis on which each rests. Four kinds of statement are
distinguished. A \emph{theorem} is proved, by an argument given in the text and by exact machine
arithmetic. A \emph{computed result} follows from the model by a numerical calculation whose
approximations are named. A \emph{hypothesis} is a statement that was examined, with the outcome
given. An \emph{assumption} is a statement on which other statements rest and which was not derived.

\begin{table*}
\caption{\label{tab:status}The principal statements of this paper, their nature and their basis.}
\begin{ruledtabular}
\begin{tabular}{>{\raggedright\arraybackslash}p{0.45\textwidth}>{\raggedright\arraybackslash}p{0.13\textwidth}>{\raggedright\arraybackslash}p{0.27\textwidth}>{\raggedright\arraybackslash}p{0.07\textwidth}}
Statement & Nature & Basis & Sec. \\
\hline
%%<1644>
The eight Dirac matrices are real and satisfy the anticommutation relations of the $4+4$ metric & theorem & exact arithmetic & \ref{sec:matrices} \\
Their scaled commutators span $\so(4,4)$ and generate, with the matrices, $\Pin(4,4)$ with four components & theorem & exact arithmetic; structure of the Clifford group & \ref{sec:pin} \\
The chirality matrix and the chiral projectors are real & theorem & exact arithmetic & \ref{sec:chirality} \\
The Dirac equation $(\gamma^A\partial_A-m)\psi=0$ is covariant, has a conserved current and splits into two coupled real 8-component equations & theorem & derivation; exact arithmetic & \ref{sec:dirac} \\
The evolution in $x4$ is unitary exactly for fields that do not depend on $x5,x6,x7$; otherwise growing modes exist & theorem & exact arithmetic & \ref{sec:quantized} \\
The physical fields do not depend on the additional times & assumption & the second stage of the program would have to justify it & \ref{sec:quantized} \\
A single field has no Lorentz-invariant pair; the pair along $x0$ is a scalar of ordinary spacetime & theorem & exact arithmetic & \ref{sec:pairs} \\
In the covariant $4+4$ theory no Vorton-Scalar with a chiral phase symmetry exists & theorem & derivation & \ref{sec:nogo} \\
In the quantized sector the Vorton-Scalar exists; its neutral current couples to the axial current with $\GVS/\sqrt2=1/2\vVS^2$ & definition and derivation & Lagrangian (\ref{eq:LVS}) & \ref{sec:vs} \\
The zero-mode sector has the internal symmetry $su(2)_A\oplus su(2)_B\oplus u(1)_F\oplus u(1)_E$, of which the $SO(4)$ is hidden & computed result & linear algebra on the matrices & \ref{sec:internal} \\
%%</1644>
%%<1671>
The eight Dirac matrices are real and satisfy the anticommutation relations of the $7+1$ metric; real matrices force the convention in which the time-like matrix squares to $+1$ & theorem & exact arithmetic; structure of the Clifford algebras & \ref{sec:matrices} \\
Their scaled commutators span $\so(7,1)$ and generate, with the matrices, $\Pin(7,1)$ with four components & theorem & exact arithmetic; structure of the Clifford group & \ref{sec:pin} \\
No real chirality matrix exists; the chiral projectors are complex & theorem & exact arithmetic & \ref{sec:chirality} \\
The Dirac equation $(i\gamma^A\partial_A-m)\psi=0$ is covariant, has a conserved positive density and splits into two coupled 8-component equations & theorem & derivation; exact arithmetic & \ref{sec:dirac} \\
A single field has no Lorentz-invariant pair; the pair along $x0$ is a scalar of ordinary spacetime & theorem & exact arithmetic & \ref{sec:pairs} \\
The Vorton-Scalar exists; its neutral current couples to the axial current with $\GVS/\sqrt2=1/2\vVS^2$ & definition and derivation & Lagrangian (\ref{eq:LVS}) & \ref{sec:vs} \\
The zero-mode sector has the internal symmetry $su(2)_A\oplus su(2)_B\oplus u(1)_F\oplus u(1)_E$, the rotations of the four extra space axes and two phases & computed result & linear algebra on the matrices & \ref{sec:internal} \\
The fields that matter do not depend on the four extra coordinates & assumption & the second stage of the program would have to justify it & \ref{sec:reduction} \\
%%</1671>
One-boson exchange attracts in the vector channels along the extra axes; the neutral current repels in the two-form channels & computed result & closed forms and matrices; static, single exchange & \ref{sec:pairing} \\
A paired state exists at any density; the vector channel has the largest gap & computed result & gap equation; two-loop truncation of the CJT action & \ref{sec:cjt} \\
The CJT potential reproduces gap, critical temperature, bound-state threshold and Ginzburg-Landau coefficients & computed result & independent algorithms & \ref{sec:checks} \\
The ground states are unitary; the hidden-axis singlet is one of a degenerate family & theorem and computed result & matrix identity (\ref{eq:quartic}); CJT potential & \ref{sec:order} \\
The hidden axis is the preferred direction (H1 in its strict form) & hypothesis; open & would need the second stage & \ref{sec:ground} \\
Dense matter is in the abnormal state; it is self-bound for $C>3\pi^2/8$ & computed result & mean field & \ref{sec:normal} \\
The fermion mass melts at a temperature of the order of $\vVS$ & computed result & Hartree approximation & \ref{sec:thermal} \\
A pair condensate in the vacuum needs $\aVS>\alpha_c\ge\pi/2$ & computed result & ladder approximation & \ref{sec:vacuum} \\
The singlet superfluid is the Higgs field (H2) & hypothesis; \SIG{no geometric basis}{unproven} & seven conditions & \ref{sec:higgs} \\
The model resolves the Hubble tension (H3) & hypothesis; not supported & simulation; published analyses & \ref{sec:hubble} \\
The superfluid explains the galactic phenomenology of dark matter (H3) & hypothesis; excluded & simulation; elementary estimates & \ref{sec:dm} \\
The dark matter consists of bound pairs or clusters of the fermions & possibility; untested & --- & \ref{sec:dm} \\
\end{tabular}
\end{ruledtabular}
\end{table*}

\subsection{The three hypotheses}

\emph{H1, the superfluid.} The affirmative part of the hypothesis is established within the
approximations stated. The fermion field has a pair of ordinary spin 0 that is a scalar of ordinary
spacetime; a scalar boson exists whose exchange binds that pair and whose neutral current selects
it; and at any finite density the paired state has a lower thermodynamic potential than the unpaired
one. Two parts of the hypothesis hold in a weaker form than its wording suggests. The pairing is
``mediated by neutral currents'' in the sense that the neutral current decides the channel; the
attraction that binds at low density comes from the radial mode of the same field
(Sec.~\ref{sec:pairing}). And the singlet along the hidden axis is a ground state, but one of a
degenerate family: \SIG{in the sector of fields that do not depend on $x0$ a hidden symmetry rotates
it into the triplet along the additional times, and which of the two lies lower once that symmetry is
broken has not been computed.}{on the flat $7+1$ spacetime the direction of the order parameter in
the space of the four extra axes is chosen spontaneously, and only the second stage of the program,
which treats $x0$ differently from the other three, could prefer the hidden axis.}

\emph{H2, the Higgs field.} \SIG{On the $4+4$ spacetime the hypothesis has no geometric basis. The
algebraic pattern of the Higgs vacuum, a $(2,2)$ multiplet whose vacuum value breaks $SO(4)$ to
$SO(3)$, is present for the zero modes, but under a symmetry that is not a transformation of the
spacetime and that no known mechanism turns into a gauge symmetry. Under the rotations of the
additional times the condensate is a singlet.}{On the $7+1$ spacetime the hypothesis is unproven.
One structural condition is met in a non-trivial way: the order parameter is a $(2,2)$ multiplet of
the rotations of the four extra space axes, and its vacuum value breaks $SO(4)$ to $SO(3)$, the
pattern of the Higgs vacuum. The dynamical conditions are not met where the model can be solved, and
two estimates, of the mass of the constituents and of the additional particles, speak against it.}
In both spacetimes a condensate in empty space, which the hypothesis needs, was found only at
couplings at which the present methods are a guide and not a proof.

\emph{H3, dark matter and the Hubble tension.} Both parts were tested by simulation, and both
answers are negative in the sense that matters. The dark radiation of the model raises the Hubble
constant, but for data that agree with the standard model the price is at least
{{GEN:dr_sig_local_min_all}} standard deviations, and the published analyses of the real data leave
room for about {{GEN:dr_act_fraction}} per cent of the discrepancy. A superfluid that could act on
galactic scales erases the small-scale structure that is observed. What remains is the possibility
that heavy bound states of the fermions are cold dark matter, which the present observations neither
support nor exclude and which makes no use of the superfluid.

\subsection{What the two spacetimes share, and where they differ}

The two papers were written from one source, and their common sections are identical. This is
possible because of a fact that was not obvious at the outset: \emph{in the sector of fields that
depend on ordinary space and time only, the two models have the same structure.} Their Hamiltonians
have the same form, with different matrices; their internal symmetry algebras are isomorphic; their
exchange amplitudes agree to the accuracy of the arithmetic; and the paired state, its
thermodynamics, the behavior of the Vorton-Scalar in matter and at finite temperature, and the
cosmology are the same. The two sets of Dirac matrices are different sets of generators of the same
algebra of real $16\times16$ matrices, related by Eq.~(\ref{eq:relation}).

They differ in everything that involves the four extra axes as directions of the spacetime.
\SIG{Here, on the $4+4$ spacetime, three of them are times. The quantum theory exists only for fields
that do not depend on them. The chirality is real, and the fully covariant theory admits no
Vorton-Scalar. The invariant form on the order parameter is indefinite, the components along the
time-like axes carry a relative phase of $90^\circ$, and the $SO(4)$ that mixes the hidden axis with
the other three is a hidden symmetry of the zero modes.}{Here, on the $7+1$ spacetime, all four are
space-like. The quantum theory is conventional. The chirality is complex, and the Vorton-Scalar
exists in the fully covariant theory. The invariant form on the order parameter is positive
definite, and the $SO(4)$ that mixes the hidden axis with the other three is the rotation group of
the extra space, a symmetry of the spacetime.} The companion paper
\SIG{\cite{Companion1671}}{\cite{Companion1644}} states the corresponding results for the other
signature. For the hypotheses the difference matters in one place, the Higgs hypothesis, where the
$7+1$ spacetime offers a geometric origin of the custodial pattern and the $4+4$ spacetime does not.

\subsection{Limits of the treatment}

Five limits apply to the paper as a whole; the limits of particular calculations were stated in the
sections concerned.

\emph{An effective theory.} The couplings of the model have negative mass dimension in eight
dimensions. The model is an effective description below some scale, and nothing here depends on its
completion; but nothing here constrains that completion either.

\emph{Flat spacetime.} Gravity is absent. The roles of the coordinates, that $x0$ is hidden and that
$x5,x6,x7$ are not observed, are put in by the restriction to fields that do not depend on them. The
second stage of the program, a dynamical metric in which three dimensions inflate, three deflate and
one curls up, is what would have to justify this restriction, decide the direction of the order
parameter, and determine whether any rotation of the extra dimensions survives as a gauge symmetry.

\emph{Truncations.} The paired state was obtained at two loops of the CJT action with a static
interaction; the behavior in matter in the mean-field approximation; the thermal transition in the
Hartree approximation; the vacuum condensate in the ladder approximation. Each is the standard first
approximation of its problem, and each is known to fail quantitatively at strong coupling.
Appendix~\ref{app:toy} shows on a solvable example how the two-loop truncation behaves.

\emph{Mock data.} The cosmological conclusions about the Hubble tension rest on simulated data of the
quality of the Planck satellite and on published analyses of the real data, not on a new analysis of
the real data.

\emph{One field.} The three hypotheses were examined for one fermion field with one mass. They need
three different mass scales (Sec.~\ref{sec:higgs}); a model with several such fields would escape
this particular objection at the price of more parameters.

\subsection{What should be done next}

Three calculations would settle questions that this paper leaves open. \SIG{First, the splitting
between the singlet along the hidden axis and the triplet along the additional times when the fields
depend on $x0$, for example on a ring of radius $R$: this decides whether the singlet of hypothesis
H1 is the ground state.}{First, the energy of the order parameter when the hidden axis is curled up
into a ring of radius $R$ and the other three extra dimensions are not: this decides whether the
singlet of hypothesis H1 is the ground state.} Second, the pair condensate in the vacuum beyond the
ladder approximation, with the corrections to the vertices and to the boson propagators and with the
full 16-component multiplet of Sec.~\ref{sec:vacuum}: this decides condition C2 of the Higgs
hypothesis. Third, the formation and the abundance of bound clusters of the fermions in the early
universe: this decides whether the one surviving possibility for the dark matter is viable.

% =====================================================================
\section{Conclusions}
\label{sec:concl}
% =====================================================================

We asked whether a 16-component Dirac field on a flat eight-dimensional spacetime of signature
\SIG{$4+4$}{$7+1$} forms a spin-singlet superfluid, and whether that state can be the Higgs field,
the dark matter, or the remedy of the Hubble tension.

The mathematical foundation is secure. The eight Dirac matrices are real signed permutation
matrices; they satisfy the anticommutation relations of the metric; their scaled commutators span
\SIG{$\so(4,4)$}{$\so(7,1)$} and generate, together with the matrices, the four components of
\SIG{$\Pin(4,4)$}{$\Pin(7,1)$}. \SIG{The chirality is real, the Dirac equation is real, and it
defines a quantum theory exactly for fields that do not depend on the three additional
times.}{No real chirality exists, the Dirac equation necessarily contains a factor $i$, and the
field is quantized in the conventional way.} Every one of these statements is verified by exact
arithmetic.

The superfluid exists. We invented the Vorton-Scalar, a complex scalar field whose vacuum value is
the origin of the fermion mass and whose neutral current couples to the axial current of the
fermions with a constant of the Fermi type. Its radial mode binds two fermions, its neutral current
selects the pairs that form a four-vector along the four extra axes, and the component of that
four-vector along the hidden coordinate $x0$ is a spin-singlet scalar of ordinary spacetime. Two
independent formulations of the many-body problem, the gap equation and the effective potential of
Cornwall, Jackiw and Tomboulis, agree on the paired state. The singlet is a ground state, one of a
degenerate family.

The three applications fare differently, and none succeeds. \SIG{The Higgs hypothesis has no
geometric basis on this spacetime.}{The Higgs hypothesis is unproven: the order parameter has the
representation and the breaking pattern of the Higgs field, but no dynamical condition is met where
the model can be solved.} The dark radiation of the model cannot raise the Hubble constant to the
locally measured value without spoiling the fit to the microwave background. A superfluid of these
fermions cannot be the dark matter of galaxies, although heavy bound states of them could be cold
dark matter of an ordinary kind.

These negative results are results. They are quantitative, they come with the programs that
produced them, and they mark the places where the program to which this paper belongs must supply
something new: a dynamical metric that distinguishes the hidden coordinate, and a mechanism that
operates at strong coupling.

\section*{Data availability}

The source code of all programs, the data files from which every figure and every computed number of
this paper is generated, the instructions for compiling and running the programs, and the documents
that record the provenance of each result are openly available \cite{SpacetimesRepo}.
